\chapter{What this version completes}

The first usable version is a terminal application for a single user. It answers questions through Ollama, manages an explicit list of saved facts and personal tasks, retrieves public web evidence, and reads or proposes changes to small code/text files in a selected workspace. A tool is a Python function that the model may request; the model itself does not open a file or save a task. The application validates the request, runs an allowed function, and sends its actual result back to the model.

The original scripts combined configuration, JSON persistence, tool schemas, networking, system instructions and an interactive loop in one file. The main script exceeded a thousand lines. The current code separates those responsibilities without introducing an agent framework, database, embeddings, vector store or server. You can follow each request through a handful of Python modules.

The concrete startup problems were a missing environment-file loader, a mismatch between the key name used by the source and the one present locally, and missing dependencies in the terminal's selected Python environment. Other weaknesses were state paths dependent on the current directory, silent resets on invalid JSON, direct overwrites, reusable task IDs after deletion, unbounded conversation growth, inconsistent tool argument validation, verbose raw exceptions and repeated prompt rules.

The existing key name works through a compatibility alias. The existing private JSON files are preserved and copied into a dedicated local state directory. The old scripts are preserved under examples/legacy and their former root filenames now launch the maintained application. No private state or key value belongs in the report or source appendix.

This release intentionally gives the user control over model-requested changes: a proposed task/memory mutation shows its arguments, and a code change shows a diff. Only an explicit YES applies it. Direct slash commands express the user's action directly, so adding a task or fact through them is immediate; deletion still requires confirmation. The application does not execute arbitrary shell commands or tests generated by a model.

The code and libraries do not require a paid subscription. Cloud inference can use available Ollama account credits, but those credits and eligible models have service limits. This is not unlimited free cloud. Local inference avoids cloud API usage, at the cost of local hardware, download bandwidth, storage and electricity. The tested default cloud model was gemma4:31b. A local server was not tested in this session.

